\section{Related Work}
\label{sec:related}
\subsection{Research Direction A}
Group relevant prior studies by the ideas they share. Explain the problem
they solve, their key assumptions, and the evidence they provide.
Start with a foundational method~\cite{foundational} and use a survey to
explain how that method fits into the wider field~\cite{survey}.
Replace these fictional entries with real papers you have read.

Describe the assumptions that make this family of methods suitable for its
target setting. Identify the information it requires and the type of output
it produces. When discussing a shortcoming, distinguish a limitation stated
by the original authors from your own interpretation. Explain the evidence
behind that interpretation and avoid extrapolating beyond the study's scope.

\subsection{Research Direction B}
Describe an alternative family of approaches and compare its strengths
and limitations with the first group. Use consistent comparison criteria,
such as required inputs, applicability, or evaluation conditions.
Use the alternative-method citation for the approach itself~\cite{alternative}
and the representation citation for the information it operates
on~\cite{representation}. This separation demonstrates how different sources
can support distinct parts of a comparison.

Compare the two directions under the same criteria. A method that uses
additional input information may not be directly comparable to one designed
for a more constrained setting. Explain such differences before discussing
performance. If the available studies use different protocols, describe
those differences rather than treating their reported values as equivalent.

\subsection{Hybrid Approaches and Shared Components}
Some papers combine ideas from multiple research directions. Use a source
on a hybrid approach to discuss that relationship~\cite{hybrid}, then identify
the components it shares with earlier work~\cite{foundational,alternative}.
Explain whether the combination changes the assumptions, the computational
requirements, or the type of problem that can be addressed.

When your paper reuses a representation or solver, acknowledge those
dependencies explicitly~\cite{representation,optimization}. Distinguish the
origin of a method from the implementation used in your experiments. A
software reference~\cite{software} credits the implementation, but it does
not replace a citation to the underlying scientific idea.

\subsection{Evaluation Practices}
Review how related studies evaluate their claims. Dataset documentation
and metric definitions serve different purposes~\cite{dataset,metric}:
the former describes the evaluation material, while the latter explains
what a reported number means. Discuss whether the protocol measures the
capability that your paper aims to improve.

Identify the most relevant baseline~\cite{baseline} and explain why it is
an informative comparator. Use guidance on repeated trials and controlled
experiments to motivate your evaluation design~\cite{reproducibility,ablation}.
Replace these placeholders with sources appropriate to your discipline;
the template does not assume that one protocol fits every research problem.

\subsection{Positioning of This Work}
Explain precisely how your question or approach relates to the reviewed
work. Identify what you reuse and what you change. A difference in method
does not by itself establish an improvement; state what evidence would
be needed to support that claim.
Relate the proposed contribution to the closest method, rather than only
to a broad survey~\cite{hybrid}. State which previously studied limitation
your evaluation will address~\cite{limitations}, and which questions remain
outside the paper. This paragraph should lead directly to the problem
formulation and the comparisons used later in the manuscript.
